\section{Corrections to an Earlier Audit (RQ4)}\label{sec:carry}
An earlier audit of this system's single-controller version (a mechanism-by-fault ablation, E1, over five registered fault classes) contained claims that did not survive review. We list each with the correction and its evidence; none of the earlier stored results was overwritten.
\paragraph{Prediction count.} The earlier manuscript reported that all 14 registered predictions held. The registered prediction for the no-guardrail variant (V2) was that invariant I7 (no difficulty increase after a score below 0.30) would be violated \emph{more often} than in the full system. Recomputing from the stored runs, I7 was violated in 3 of 8 sessions for V2 and 7 of 8 for the full system, so \textbf{13 of 14 predictions held and one was contradicted}; the earlier analysis had applied a weaker criterion (any violation).
\paragraph{Malformed scorer output.} The earlier text implied the shipped system committed malformed scores in production. At the hub seam with an injected evaluator (synthetic fault injection), a non-numeric, \texttt{None}, NaN, out-of-range or string score was committed after repair to 0.0 or a clamped value in every session of that cell. In the deployed backend wrapper, however, \texttt{None} and non-numeric scores already raise (so the hub fails closed), and only NaN (to 0.0) and out-of-range numbers (1.7 to 1.0, $-0.3$ to 0.0) are repaired; the evaluator itself clamps to $[0,1]$. A regression test pins this behaviour. The strict contract is therefore defence in depth at the hub seam, not the fix of a demonstrated production defect.
\paragraph{Historical reproducibility.} The strict contract is now default-on, so re-running E1 with defaults gives different results in the malformed-score cells. Under explicit historical flags (strict contract off, breaker disabled) a re-run of all \xeEoneRows{} sessions was byte-identical to the stored file (\xeEoneDiff{} sessions differing; SHA-256 equal). Experiments therefore carry an explicit configuration.
\paragraph{Guardrail divergence.} The deployed orchestrator uses an inline rule layer that differs from the module used by the training environment and the earlier policy study. An earlier figure of 12.4\% disagreement on a synthetic grid included an artefact of the comparison script (the canonical function reads difficulty values at most 1.0 as already normalised); with a normalised input the disagreement is 4,728 of 39,930 points (11.8\%; 7.4\% on a full grid). A configuration key selects either implementation and both are tested; the shipped default is unchanged.
\paragraph{Scorer semantics.} The evaluator's calibration assumes a sigmoid-bounded cross-encoder score. The pinned library range applies a sigmoid by default for a single-label head (verified in the source of version 2.7.0), but the installed 6.1.0 returns raw logits (values 1.55 and $-7.95$ for two test pairs, equal to the sigmoid-inverse of the explicit-activation outputs). The loader now passes the activation explicitly; a test checks the loader for both keyword conventions and, where weights are materialised, that scores are probabilities equal to the sigmoid of the logits.
\paragraph{Grounding artefact.} See Sect.~\ref{sec:results}: the 17 false rejections were artefacts, and the unsupported-credit acceptance is higher than reported.
\paragraph{Circuit breaker.} In the earlier breaker, any HTTP response counted as success. Against a service that answers every request with HTTP 503, the historical semantics sent \bfivezerothreeerrfivezerothreelegacy{} requests over a 24-turn replay session at a median \bfivezerothreemederrfivezerothreelegacy{}\,s per turn, whereas treating 5xx as failure sent \bfivezerothreeerrfivezerothreefixed{} requests at \bfivezerothreemederrfivezerothreefixed{}\,s per turn; the 24 committed scores were identical in all variants (SOCKET, one session). A cancelled half-open probe previously left the breaker rejecting all calls; it is now released.
\paragraph{Breaker timing re-run.} The breaker timings were re-measured with current code against a local socket server with real six-second timeouts (a new SOCKET run, not a reproduction; one 24-turn replay session). Per-turn medians in seconds, new (earlier): hanging service without the breaker \efivehangBnew{} (\efivehangBold{}); with the breaker at pacing $P{=}30$, 60 and 120\,s \efivehangSPthreezeronew{} (\efivehangSPthreezeroold{}), \efivehangSPsixzeronew{} (\efivehangSPsixzeroold{}) and \efivehangSPonetwozeronew{} (\efivehangSPonetwozeroold{}); healthy \efiveokSPsixzeronew{} (\efiveokSPsixzeroold{}), refusing \efiverefuseSPsixzeronew{} (\efiverefuseSPsixzeroold{}), recovery \efiverecoverSPsixzeronew{} (\efiverecoverSPsixzeroold{}). HTTP attempts were unchanged in the hanging scenarios (\efivehangBnewhttp{}, \efivehangSPthreezeronewhttp{}, \efivehangSPsixzeronewhttp{}). One run on one machine, so differences of a few hundredths of a second are not interpreted.
